\documentclass[UTF8,a4paper,11pt,openany]{ctexbook}
\usepackage{../common/statlearnbook}
\SLSetBookMetadata{第 1 册}{统计学习方法讲义 v2.6.0}{数学基础与统计学习基本理论}
\begin{document}
\setcounter{page}{202}
\setcounter{chapter}{11}
\setcounter{figure}{0}
\pagestyle{headings}
\chaptermark{监督学习任务与应用}
\begin{center}
\begin{minipage}[t]{0.47\linewidth}
\begin{keypointbox}[title=HMM任务卡]
\textbf{输入：}观测序列$X_{1:n}$。\par
\textbf{输出：}最可能的隐状态路径$Y_{1:n}$。\par
\textbf{箭头语义：}状态生成下一状态与当前观测，属于生成式联合模型。
\end{keypointbox}
\end{minipage}\hfill
\begin{minipage}[t]{0.47\linewidth}
\begin{keypointbox}[title=CRF任务卡]
\textbf{输入：}完整观测序列$x_{1:n}$。\par
\textbf{输出：}条件得分最高的标记路径$y_{1:n}$。\par
\textbf{因子语义：}局部因子在给定观测下连接相邻标记，不表示由标记生成观测。
\end{keypointbox}
\end{minipage}
\end{center}
% v2.3.1 figure UID: FIG-P172-01
% Proposed post-v2.3.1 polish for FIG-P172-01: protect key-label size and stroke hierarchy.
\tikzset{slfig-FIG-P172-01/.style={font=\fontsize{9.2pt}{11.0pt}\selectfont,
  every path/.append style={line cap=round,line join=round}}}
\begin{figure}[H]
\centering
\begin{tikzpicture}[slfig-FIG-P172-01,
  paneltitle/.style={font=\bfseries\fontsize{10pt}{12pt}\selectfont,align=center},
  latent/.style={slfig latent,minimum size=7.5mm,line width=.74pt},
  observed/.style={slfig observed,minimum size=7.5mm,line width=.74pt,
    fill=SLBlueBG,pattern={Lines[angle=45,distance=4.4pt,line width=.18pt]},
    pattern color=SLBlue!14},
  factor/.style={slfig factor,fill=SLGoldBG,draw=SLGold,line width=.78pt},
  gen/.style={-{Stealth[length=2.05mm,width=1.25mm]},draw=SLBlue,
    line width=.94pt,shorten >=.55mm},
  undirected/.style={draw=SLInk,line width=.78pt},
  conditionbrace/.style={draw=SLTeal,line width=.64pt},
  legendlabel/.style={font=\fontsize{9.2pt}{11pt}\selectfont,anchor=west},
]
  \begin{scope}[xshift=-5.2cm]
    \node[paneltitle,text=SLBlue] at (2.4,2.45) {HMM：生成关系};
    \node[latent] (hy1) at (0,1.25) {$Y_t$};
    \node[latent] (hy2) at (1.45,1.25) {$Y_{t+1}$};
    \node at (2.88,1.25) {$\cdots$};
    \node[latent] (hy3) at (4.35,1.25) {$Y_T$};
    \node[observed] (hx1) at (0,0) {$X_t$};
    \node[observed] (hx2) at (1.45,0) {$X_{t+1}$};
    \node at (2.88,0) {$\cdots$};
    \node[observed] (hx3) at (4.35,0) {$X_T$};
    \draw[gen] (hy1)--(hy2);
    \draw[gen] (hy2)--(2.68,1.25);
    \draw[gen] (3.08,1.25)--(hy3);
    \draw[gen] (hy1)--(hx1);
    \draw[gen] (hy2)--(hx2);
    \draw[gen] (hy3)--(hx3);
  \end{scope}

  \begin{scope}[xshift=1.0cm]
    \node[paneltitle,text=SLTeal] at (2.4,2.45) {线性链 CRF：条件因子};
    \node[latent] (cy1) at (0,1.25) {$Y_t$};
    \node[latent] (cy2) at (1.45,1.25) {$Y_{t+1}$};
    \node at (2.88,1.25) {$\cdots$};
    \node[latent] (cy3) at (4.35,1.25) {$Y_T$};
    \node[observed] (cx1) at (0,0) {$X_t$};
    \node[observed] (cx2) at (1.45,0) {$X_{t+1}$};
    \node at (2.88,0) {$\cdots$};
    \node[observed] (cx3) at (4.35,0) {$X_T$};
    \node[factor] (f12) at (.725,1.25) {};
    \node[factor] (f23a) at (2.62,1.25) {};
    \node[factor] (f23b) at (3.14,1.25) {};
    \node[factor] (g1) at (0,.625) {};
    \node[factor] (g2) at (1.45,.625) {};
    \node[factor] (g3) at (4.35,.625) {};
    \draw[undirected] (cy1)--(f12)--(cy2);
    \draw[undirected] (cy2)--(f23a);
    \draw[undirected] (f23b)--(cy3);
    \draw[undirected] (cy1)--(g1)--(cx1);
    \draw[undirected] (cy2)--(g2)--(cx2);
    \draw[undirected] (cy3)--(g3)--(cx3);
    \draw[decorate,decoration={brace,mirror,amplitude=3pt},conditionbrace]
      (-.35,-.45)--(4.70,-.45)
      node[midway,below=2pt,font=\fontsize{9.2pt}{11pt}\selectfont,text=SLTeal,
        fill=white,inner sep=.4pt] {给定观测 $\boldsymbol x$};
  \end{scope}

  \node[latent,minimum size=5.8mm] at (-3.95,-1.15) {};
  \node[legendlabel] at (-3.60,-1.15) {隐变量};
  \node[observed,minimum size=5.8mm] at (-2.15,-1.15) {};
  \node[legendlabel] at (-1.80,-1.15) {观测变量};
  \node[factor,minimum size=3.4mm] at (.15,-1.15) {};
  \node[legendlabel] at (.42,-1.15) {条件因子（无向）};
  \draw[gen] (2.60,-1.15)--(3.20,-1.15);
  \node[legendlabel] at (3.28,-1.15) {生成方向};
\end{tikzpicture}
\caption{HMM与线性链CRF都利用相邻标记关系，但箭头只用于HMM的生成方向；CRF中的无向线段表示给定观测后的局部因子。}\label{fig:V1-C11-tagging}
\end{figure}

由\cref{fig:V1-C11-tagging}可见，改变一个$y_t$会同时影响本地观测项和相邻转移项；逐位置取最大值因此可能拼出全局低分序列，解码必须保留跨位置的动态规划状态。
\end{document}
